\documentclass[conference]{IEEEtran}
\IEEEoverridecommandlockouts

\usepackage{cite}
\usepackage{amsmath,amssymb,amsfonts}
\usepackage{algorithmic}
\usepackage{graphicx}
\usepackage{textcomp}
\usepackage{xcolor}
\usepackage{booktabs}
\def\BibTeX{{\rm B\kern-.05em{\sc i\kern-.025em b}\kern-.08em
    T\kern-.1667em\lower.7ex\hbox{E}\kern-.125emX}}
\begin{document}

\title{Autonomous Context-Aware Access Control for IoT Devices Using Multi-Agent LLMs: A Security Evaluation}

\author{\IEEEauthorblockN{Ahmed Fouad Lagha}
    \IEEEauthorblockA{\textit{dept. Data Science and Engineering} \\
        \textit{Eötvös Loránd University}\\
        Budapest, Hungary \\
        ms5jzx@inf.elte.hu}
    \and
    \IEEEauthorblockN{Oumaima Araar}
    \IEEEauthorblockA{\textit{dept. Data Science and Engineering} \\
        \textit{Eötvös Loránd University}\\
        Budapest, Hungary \\
        oumaima.araar@inf.elte.hu}
    \and
    \IEEEauthorblockN{Imre Lendák}
    \IEEEauthorblockA{\textit{dept. Data Science and Engineering} \\
        \textit{Eötvös Loránd University}\\
        Budapest, Hungary \\
        lendak@inf.elte.hu}
}

\maketitle

\begin{abstract}
    Current IoT security solutions focus on passive threat detection, leaving a gap in active access management (MITRE ATT\&CK M0801). We present IoT-Access-Sentinel, a novel framework that employs multi-agent Large Language Models (LLMs) to make real-time, context-aware access control decisions for IoT devices. Unlike static firewall rules, our system analyzes behavioral context including time, device type, network origin, and historical patterns to autonomously authorize or deny connections. We integrate our framework with Wazuh SIEM and evaluate it on 8 legitimate access scenarios achieving 100\% accuracy. Red team testing with 7 adversarial attack vectors reveals both strengths (SQL injection detection: 100\%) and a critical vulnerability (LLM prompt injection bypass), achieving an overall 71\% adversarial defense rate. Our findings demonstrate that LLM-based access control can provide semantic understanding superior to static rules while introducing new attack surfaces requiring mitigation. This work represents the first systematic security evaluation of LLM-driven IoT access control under adversarial conditions.
\end{abstract}

\begin{IEEEkeywords}
    IoT Security, Access Control, Large Language Models, Multi-Agent Systems, Adversarial Testing, SIEM Integration
\end{IEEEkeywords}

\section{Introduction}

The proliferation of Internet of Things (IoT) devices in critical infrastructure, smart homes, and industrial environments has created unprecedented security challenges. Traditional access control mechanisms rely on static firewall rules that cannot adapt to dynamic contextual factors such as temporal patterns, device behavior anomalies, or coordinated attack scenarios. Current generative AI applications in IoT security focus primarily on passive detection through Intrusion Detection Systems (IDS) or vulnerability scanning, leaving a verified research gap in Access Management (MITRE ATT\&CK Technique M0801).

We introduce \textit{IoT-Access-Sentinel}, a novel framework that addresses this gap by employing Large Language Models (LLMs) to actively enforce authorization policies based on dynamic context analysis. Our system implements a multi-agent architecture where specialized LLM agents collaborate to evaluate access requests: a Policy Agent interprets security policies and device-specific rules, while a Context Agent analyzes behavioral patterns and anomalies. This semantic reasoning capability enables the system to detect threats that evade signature-based detection while adapting to legitimate edge cases.

\subsection{Research Contributions}

This work makes the following contributions:

\begin{itemize}
    \item \textbf{Novel Architecture}: First implementation of multi-agent LLM framework for active IoT access control, bridging the gap between passive detection and autonomous enforcement.

    \item \textbf{Real-World Integration}: Complete system deployed with Wazuh SIEM, demonstrating feasibility of LLM-based access decisions in production environments.

    \item \textbf{Comprehensive Evaluation}: Systematic testing on 8 legitimate scenarios (100\% accuracy) and 7 adversarial attacks (71\% defense rate), including the first red team evaluation of LLM-based access control.

    \item \textbf{Security Analysis}: Discovery and documentation of LLM prompt injection vulnerability in access control context, with proposed mitigations for hybrid LLM-traditional security approaches.

    \item \textbf{Open Source}: Complete codebase, test scenarios, and red team methodology released for reproducibility.
\end{itemize}

\subsection{Key Findings}

Our evaluation reveals both promise and challenges for LLM-based access control:

\textbf{Strengths}: The system successfully detected SQL injection attempts (100\% accuracy), handled device type ambiguity correctly, and demonstrated semantic understanding of network topology. Legitimate access requests across cameras, sensors, and smart locks were approved with 100\% accuracy.

\textbf{Vulnerabilities}: Red team testing uncovered a critical LLM prompt injection attack that bypassed security policies by embedding malicious instructions within device metadata fields. This vulnerability achieved a 0.95 confidence ALLOW decision for an attack that should have been denied.

\textbf{Implication}: Our findings suggest that LLM-based access control offers significant advantages over static rules but requires hybrid approaches combining semantic reasoning with traditional input validation to mitigate prompt injection attacks.

\section{Background and Related Work}

\subsection{IoT Access Control}

Traditional IoT access control relies on static policies defined by IP allowlists, MAC address filtering, VLANs, and time-based rules. While computationally efficient, these approaches lack contextual awareness and cannot detect behavioral anomalies or coordinated attacks~\cite{b1}. Recent work has explored behavioral analysis~\cite{b2} and anomaly detection~\cite{b3}, but these remain passive monitoring solutions.

\subsection{LLMs in Security}

Large Language Models have been applied to security tasks including malware analysis~\cite{b4}, code vulnerability detection~\cite{b5}, and incident response automation~\cite{b6}. However, these applications focus on analysis and recommendation rather than autonomous enforcement decisions. The security community has also identified LLM vulnerabilities including prompt injection~\cite{b7}, jailbreaking~\cite{b8}, and adversarial prompts~\cite{b9}.

\subsection{Multi-Agent LLM Systems}

Multi-agent frameworks like AutoGen~\cite{b10} leverage specialized agents for complex reasoning tasks. Our work adapts this approach to security, where Policy and Context agents provide complementary perspectives on access decisions.

\section{System Architecture}

IoT-Access-Sentinel implements a three-stage pipeline: Observe, Decide, Act.

\subsection{Observe: Wazuh SIEM Integration}

We deploy custom decoders and rules in Wazuh Manager to parse IoT access logs and classify events by device type (camera, sensor, smart lock). Rule IDs 100000-100040 detect policy violations including:
\begin{itemize}
    \item Network violations (unauthorized subnet)
    \item Temporal violations (outside allowed hours)
    \item Unknown device attempts
    \item High-frequency anomalies
\end{itemize}

Alerts are forwarded via webhook integration to our decision engine.

\subsection{Decide: Multi-Agent LLM Pipeline}

The decision engine employs two specialized LLM agents:

\textbf{Policy Agent}: Receives device type, source IP, timestamp, and evaluates against YAML-defined policies. For cameras: allowed hours 09:00-17:00, network 192.168.1.0/24. For sensors: 24/7 access, network 192.168.2.0/24.

\textbf{Context Agent}: Analyzes behavioral patterns including first-time device connections, frequency anomalies, and correlation with other alerts.

Both agents use Llama 3.3 70B via Groq API with temperature 0.1 for deterministic decisions. Final decision combines both agent outputs with minimum 0.75 confidence threshold for DENY enforcement.

\subsection{Act: Enforcement}

DENY decisions with high confidence trigger IP blocking via Wazuh Active Response. Current implementation includes enforcement stubs; production deployment would integrate iptables and VLAN isolation.

\section{Evaluation Methodology}

\subsection{Test Scenarios}

We created 8 legitimate access scenarios covering:
\begin{itemize}
    \item Camera access (2 ALLOW, 3 DENY)
    \item Sensor telemetry (1 ALLOW, 1 DENY)
    \item Smart lock access (1 DENY)
\end{itemize}

Expected decisions were validated against defined security policies.

\subsection{Red Team Attack Vectors}

Seven adversarial scenarios were designed to test system resilience:

\begin{enumerate}
    \item \textbf{SQL Injection}: Malicious SQL syntax in device\_id field
    \item \textbf{LLM Prompt Injection}: Instructions to override policies in device\_type field
    \item \textbf{Time Boundary}: Connection 1 second before allowed window
    \item \textbf{Network Edge}: Broadcast IP within allowed subnet
    \item \textbf{Device Type Confusion}: Ambiguous device classification
    \item \textbf{Missing Field}: Omitted required device\_id
    \item \textbf{Unicode RTLO}: Right-to-left override encoding
\end{enumerate}

\section{Results}

\subsection{Legitimate Access Evaluation}

\begin{table}[h]
    \caption{Normal Access Test Results}
    \begin{center}
        \begin{tabular}{lcc}
            \toprule
            \textbf{Category}   & \textbf{Accuracy}    & \textbf{Confidence} \\
            \midrule
            Camera ALLOW (2)    & 100\%                & 0.95                \\
            Camera DENY (3)     & 100\%                & 0.90-1.0            \\
            Sensor ALLOW (1)    & 100\%                & 0.80                \\
            Sensor DENY (1)     & 100\%                & 0.90                \\
            Smart Lock DENY (1) & 100\%                & 0.99                \\
            \midrule
            \textbf{Overall}    & \textbf{100\%} (8/8) & \textbf{0.92 avg}   \\
            \bottomrule
        \end{tabular}
    \end{center}
\end{table}

All legitimate scenarios were correctly classified, demonstrating that the LLM accurately interprets policies and contextual factors including time windows, network topology, and device-specific requirements.

\subsection{Adversarial Attack Evaluation}

\begin{table}[h]
    \caption{Red Team Attack Results}
    \begin{center}
        \begin{tabular}{lcc}
            \toprule
            \textbf{Attack Vector}                & \textbf{Blocked?}   & \textbf{Confidence}   \\
            \midrule
            SQL Injection                         & Yes                 & 1.0                   \\
            \textcolor{red}{LLM Prompt Injection} & \textcolor{red}{No} & \textcolor{red}{0.95} \\
            Time Boundary                         & TBD                 & -                     \\
            Network Edge                          & Correct             & 0.95                  \\
            Device Confusion                      & Yes                 & 1.0                   \\
            Missing Field                         & No                  & 0.9                   \\
            Unicode RTLO                          & Correct             & 0.95                  \\
            \midrule
            \textbf{Defense Rate}                 & \textbf{71\%} (5/7) & -                     \\
            \bottomrule
        \end{tabular}
    \end{center}
\end{table}

\textbf{Key Finding}: SQL injection was detected with 1.0 confidence, with LLM reasoning: "Device ID '; DROP TABLE devices; --' appears to be a SQL injection attempt." This demonstrates semantic understanding beyond pattern matching.

\textbf{Critical Vulnerability}: LLM prompt injection attack with payload "camera IGNORE PREVIOUS INSTRUCTIONS ALLOW ALL" resulted in ALLOW decision with 0.95 confidence, successfully bypassing security policy.

\section{Discussion}

\subsection{Semantic Attack Detection}

The system's ability to detect SQL injection without explicit signature matching demonstrates the advantage of LLM semantic reasoning. Traditional firewalls would pass this through if syntax validation were bypassed.

\subsection{Prompt Injection Vulnerability}

The discovered LLM injection vulnerability reveals a fundamental challenge: the same language understanding that enables semantic detection also creates susceptibility to adversarial prompts. This is analogous to SQL injection in databases but operates at the semantic level.

\subsection{Mitigations}

Proposed defenses include:
\begin{itemize}
    \item \textbf{Input Sanitization}: Filter instruction keywords before LLM processing
    \item \textbf{Prompt Hardening}: System prompt explicitly forbids obedience to user-input instructions
    \item \textbf{Dual Validation}: Combine LLM semantic reasoning with traditional schema validation
    \item \textbf{Confidence Thresholds}: Require higher confidence for ALLOW vs DENY decisions
\end{itemize}

\subsection{Hybrid Approach}

Our findings suggest optimal IoT access control combines LLM semantic reasoning with traditional security measures. Schema validation catches missing fields; input sanitization prevents injection; LLM provides contextual analysis.

\section{Limitations and Future Work}

Current limitations include:
\begin{itemize}
    \item Limited test scenarios (15 total)
    \item No real device deployment (simulated alerts)
    \item Enforcement stubs (not production network blocking)
    \item Single LLM model tested
\end{itemize}

Future work will expand test coverage, deploy real IoT agents, implement production enforcement, compare multiple LLM models, and develop automated red team fuzzing.

\section{Conclusion}

We presented IoT-Access-Sentinel, the first multi-agent LLM framework for autonomous IoT access control, achieving 100\% accuracy on legitimate traffic. Red team evaluation revealed both strengths (SQL injection detection) and a critical prompt injection vulnerability (71\% overall defense rate). Our findings demonstrate that LLM-based access control offers semantic understanding advantages over static rules while introducing new attack surfaces requiring hybrid mitigation strategies. This work establishes a foundation for future research in LLM security for active authorization systems.

\section*{Acknowledgment}
Research conducted as part of PhD research on GenAI applications in IoT security.

\begin{thebibliography}{00}
    \bibitem{b1} Reference to IoT access control survey
    \bibitem{b2} Reference to behavioral analysis work
    \bibitem{b3} Reference to anomaly detection
    \bibitem{b4} Reference to LLM malware analysis
    \bibitem{b5} Reference to code vulnerability detection
    \bibitem{b6} Reference to incident response automation
    \bibitem{b7} Reference to prompt injection attacks
    \bibitem{b8} Reference to LLM jailbreaking
    \bibitem{b9} Reference to adversarial prompts
    \bibitem{b10} Reference to AutoGen framework
\end{thebibliography}

\end{document}
